\SourceLeaf{280}
\begin{list}{}{\leftmargin=2em\rightmargin=0pt\itemindent=0pt\itemsep=0pt\parsep=0pt}
\item erstens die Dreiecke sich nicht gar zu weit von dem Normal-Parallelkreise entfernen, und
\item zweitens, dass sie vergleichungsweise, nemlich nach dem Verhältnisse der Seiten zu einem ganzen Erdquadranten, klein sind, wie bei wirklich messbaren Dreiecken immer der Fall ist.
\end{list}
\noindent\hspace*{\parindent}Dieses genaue Anschmiegen der auf die Kugelfläche übertragenen Dreiecksseiten an Grösste Kreisbögen findet nun bei der in Obigem betrachteten conformen Darstellung in noch viel höherm Grade Statt, als bei der a.\,a.\,O. vorgeschlagenen. Wo diese [nach Art. 13] bei einem Abstände von \(2\tfrac{1}{2}\) Grad von dem Normal-Parallelkreise eine linearische Vergrösserung von \(\tfrac{1}{530000}\) ergab, würde die neue Methode nur eine Aenderung von \(\tfrac{1}{5800000}\) geben.

Man kann daher das ganze System, nachdem man zuvörderst eine Dreiecksseite auf die Kugelfläche gehörig übertragen hat, ganz so, als wenn es auf dieser selbst läge, vermittelst der Winkel berechnen, nöthigenfalls mit der eben angedeuteten Modification, sodann für alle Punkte die Werthe der Breiten und Längen bestimmen, und von diesen vermittelst der oben gegebenen Formeln, oder vielmehr was die Breiten betrifft, vermittelst einer solchen Hülfstafel, wie hier beigefügt ist, auf die Breiten und Längen auf der Ellipsoidfläche übergehen.
\Article{12}
\indent Es bleibt demnach hier noch übrig, die Bestimmung der Abweichung einer auf die Kugelfläche übertragenen geodaetischen Linie von dem zwischen denselben Endpunkten enthaltenen Grössten Kreisbogen zu entwickeln, wonach sich zugleich in jedem Falle beurtheilen lässt, ob die Berücksichtigung dieser Abweichung nöthig werde. Man kann diese Aufgabe auf mehr als eine Art behandeln: für den gegenwärtigen Zweck, wo die Reduction immer nur eine sehr kleine Grösse betragen kann, scheint folgende Methode die angemessenste zu sein.

Es sei \(L\) die in Rede stehende geodaetische Linie auf dem Ellipsoid in unbestimmter Ausdehnung betrachtet, \(M\) ihre conforme Darstellung auf der Kugelfläche, \(F\) und \(G\) die Endpunkte eines bestimmten Stückes von \(M\), endlich \(N\) ein durch diese beiden Punkte geführter Grösster Kreis. Jeder Punkt in \(N\) werde bestimmt durch seinen Abstand \(x\) von einem zunächst willkürlich auf \(N\) gewählten Anfangspunkte; jeder Punkt von \(M\) durch seinen senkrechten Abstand \(y\) von \(N\) und durch das dem Fusspunkte dieses Perpendikels zukommende
